% fix spacing
\RedeclareSectionCommand[beforeskip=-2.8ex plus -0.8ex minus -0.2ex, afterskip=1.2ex plus 0.2ex]{section}
\RedeclareSectionCommand[beforeskip=-2.6ex plus -0.8ex minus -0.2ex, afterskip=1.0ex plus 0.2ex]{subsection}
\RedeclareSectionCommand[beforeskip=-2.6ex plus -0.8ex minus -0.2ex, afterskip=1.0ex plus 0.2ex]{subsubsection}
\tcbset{btbox/.append style={before skip=0.5\baselineskip, after skip=0.5\baselineskip},
        blackbox/.style={btsidebar, btbg=ThmBlue!5, borderline west={2.6pt}{0pt}{ThmBlue}}}

% recoloring stuff
\definecolor{ThmBlue}{RGB}{155,27,48}
\definecolor{HeadSlate}{RGB}{50,57,66}
\definecolor{NumAccent}{RGB}{0,118,124}
\definecolor{DefDark}{RGB}{0,84,90}
\definecolor{LinkRose}{RGB}{0,100,108}
\colorlet{DefGreen}{NumAccent}\colorlet{ExOrange}{NumAccent}\colorlet{ParaOrange}{ThmBlue}

\renewcommand*{\subsectionformat}{{\color{NumAccent}\thesubsection}\enskip}
\renewcommand*{\subsubsectionformat}{{\color{NumAccent}\thesubsubsection}\enskip}

% unbold subtitle
\setkomafont{subtitle}{\normalfont\large\color{HeadSlate!85}}

% units
\newcommand{\unit}[1]{\,\text{\small #1}}

% this genuinely took forever to get right
\newcommand{\answerbox}[1]{\tcbox[on line, colback=ThmBlue!9, colframe=ThmBlue!65,  boxrule=0.6pt, arc=2pt, left=4pt, right=4pt, top=2pt, bottom=2pt]{$\displaystyle #1$}}
% now imagine trying to use tikz
% oh wait thats next deliverable
% kill me now

\newcommand{\journalname}{Journal of N$\Psi$ Bridge Design}
\subject{\journalname}
\setkomafont{subject}{\normalfont\large\scshape\color{HeadSlate!85}}
\title{Design Report}
\renewcommand{\thetitle}{\journalname: Design Report}

\subtitle{Team 809 \quad$\star$\quad TA: Mohi}
\author{Daniel Kessler, Cary Li, Zekai Liu, Brendan Tang, Leon Zhang}
\date{\today}

\addtokomafont{pageheadfoot}{\color{black!60}}
\addtokomafont{headsepline}{\color{ThmBlue!40}}

% get last page
\usepackage{lastpage}

\cfoot*{}\ohead{\textit{Journal of N$\Psi$ Bridge Design: Design Report}\quad\pagemark\,/\,\pageref*{LastPage}}

% GRADIENTS
% THANKS SO MUCH LATEX STACK EXCHANGE I BARELY UNDERSTAND THIS LOL
\usetikzlibrary{fadings}
\definecolor{GradGreen}{RGB}{196,150,12} % green except its not
\pgfdeclarehorizontalshading{btgrad}{100bp}{color(0bp)=(GradGreen); color(38bp)=(GradGreen); color(62bp)=(NumAccent); color(100bp)=(NumAccent)}
\newcounter{gradtext}
\newcommand{\gradtext}[1]{%
  \stepcounter{gradtext}%
  \begin{tikzfadingfrompicture}[name=gradfade\thegradtext]
    \node[inner sep=0pt, text=transparent!0] {#1};
  \end{tikzfadingfrompicture}%
  \tikz[baseline=(b.base)]{%
    \node[inner sep=0pt, text opacity=0] (b) {#1};
    \shade[path fading=gradfade\thegradtext, fit fading=false,
           fading transform={shift={(b.center)}}, shading=btgrad, shading angle=-45]
      (b.south west) rectangle (b.north east);}%
}
\renewcommand*{\sectionformat}{\gradtext{\numbigger[1.5]{\thesection}}\enskip}

% thanks again tex exchange 
% colors are from one light theme for vscode
\usepackage{minted}
\colorlet{CodeBg}{HeadSlate}
\setminted{linenos, xleftmargin=1.4em, numbersep=8pt, fontsize=\small, breaklines, autogobble, style=one-dark}
\renewcommand{\theFancyVerbLine}{\sffamily\scriptsize\color[RGB]{157,157,159}\arabic{FancyVerbLine}}
\AddToHook{env/minted/before}{\begin{tcolorbox}[btbox, colback=CodeBg!5, colframe=CodeBg!25, boxrule=0.5pt, arc=2pt, left=4pt, top=2pt, bottom=2pt]}
\AddToHook{env/minted/after}{\end{tcolorbox}}
\NewCommandCopy{\btOrigTextcolor}{\textcolor}
\newcommand{\btOneLight}[2]{\expandafter\def\csname btOL:#1\endcsname{#2}}
\btOneLight{0.67,0.70,0.75}{56,58,66}\btOneLight{0.78,0.47,0.87}{166,38,164}
\btOneLight{0.90,0.75,0.48}{193,132,1}\btOneLight{0.88,0.42,0.46}{228,86,73}
\btOneLight{0.38,0.69,0.94}{64,120,242}\btOneLight{0.34,0.71,0.76}{1,132,188}
\btOneLight{0.60,0.76,0.47}{80,161,79}\btOneLight{0.82,0.60,0.40}{152,104,1}
\btOneLight{0.50,0.52,0.56}{160,161,167}
\AddToHook{env/minted/begin}{\RenewDocumentCommand{\textcolor}{o m m}{%
  \ifcsname btOL:#2\endcsname\edef\btOLc{\csname btOL:#2\endcsname}\btOrigTextcolor[RGB]{\btOLc}{#3}%
  \else\IfValueTF{#1}{\btOrigTextcolor[#1]{#2}{#3}}{\btOrigTextcolor{#2}{#3}}\fi}}